
 Loss-Function: mse, mse\_norm und softmax? Sind differnzierbar -> benötigt für backprop. Argmax dient zur identifizierung, nicht für Training geeignet

 Warum nicht die mean squared cost function?\cite{hughesTrainingPhotonicNeural2018}
 $ L = 1/2 (X_L-T)^\dagger (X_L-T) $
 Ist ja fast das gleiche wie die mse oder nicht
 - Ich bin mir nicht sicher wie wir unser Target E-Feld definieren müssten (da bräuchten wir ja auch die Phase), deshalb hätte ich's mit Intensität gemacht (die im Paper machen das ja mit dem E-Feld, oder?)
 Soweit ich das gesehen/verstanden habe, haben sie auch einfach Vektoren mit 1 und sonst 0 für T benutzt
 - Aber das X-Feld ist dann die Intensität?
 Sie bezeichnen es meist als modal amplitude, also hätte ich Feld gedacht
 Aber es macht ja gar kein Sinn die Phase festzulegen, weil wir die ja gar nicht messen können. Und wenn wir das Feld festlegen tun wir das ja

 Oder meinst du eine nichtlineare Funktion für zwischen den Ebenen?
- Ne, die haben wir gar nicht. Das ist auch beschrieben in einem Paper... Das "Deep Learning with Coherent Nanophotonic Circuits Paper" \cite{shenDeepLearningCoherent2017} geht da etwas drauf ein, nur falls wir darauf uns nochmal beziehen wollen - Begriff der OUI vs ONU
Hier haben sie das auch in ihrer Herleitung berücksichtigt gehabt: \cite{Hughes2018-optica}
- Können wir die non-linear function einfach = id setzen?
Außer für den aller letzten schritt, also f\_l = id außer für f\_100 -> Und dann die Messung als Intensität als f\_100 einbauen?
Jap, das ist das einzige, was wir als non-linearity haben
Ok, das hatte ich auch überlegt

 Argumente für Normierung der Energie: Führt dazu, dass über Gesamthelligkeit anstatt räumliche Struktur klassifiziert wird, das ist dann anfällig für losses / laser-intensitätsschwankungen,...

